%% kapa: exam=1 round=22 year=2011 period=2 form=A subject=appraisal_law subject_ko=부동산관계법규 seq=35 number=75
%% source: Q-Net 감정평가사 기출문제 (cst00309) artlSeq=5036790 file=extracted/1008/5036790/2011년 제22회 감정평가사 1차시험 기출문제/2교시.hwp
%% conversion: kordoc 4.24.51 -> tools/gen1.py (auto); review=proofread-gpt-astra; flags=none
\begin{question}{75}
측량․수로조사 및 지적에 관한 법령상 지목에 관한 설명으로 옳지 않은 것은?
\begin{choices}
  \item 물을 상시적으로 직접 이용하여 벼․연(蓮)․미나리 등의 식물을 주로 재배하는 토지는 ‘답’이다.
  \item 자연의 유수(流水)가 있거나 있을 것으로 예상되는 토지는 ‘유지(溜池)’이다.
  \item 종교용지인 토지에 있는 유적․고적․기념물 등을 보호하기 위하여 구획된 토지는 ‘사적지’가 아니다.
  \item 육상에 인공으로 조성된 수산생물의 번식 또는 양식을 위한 시설을 갖춘 부지와 이에 접속된 부속시설물의 부지는 ‘양어장’이다.
  \item 산림 및 원야(原野)를 이루고 있는 수림지(樹林地)․죽림지 등의 토지는 ‘임야’이다.
\end{choices}
\end{question}
\answer{2}
